\section{Conclusion}
\label{sec:conclusion}

Partisan Bias provides what no other metric in the L-series provides: a direct,
seat-count-translatable measure of partisan asymmetry at the electoral pivot that
communicates intuitively to judicial audiences without technical explanation.
The statement ``if both parties received 50\% of the vote, Republicans would win
one additional seat'' is the kind of concrete, counterfactual claim that courts
can weigh against state constitutional standards without requiring expertise in
wasted-vote algebra or distributional statistics.

The core findings are robust. \textsc{Bisect} maps produce $|\text{Bias}| < 0.02$
across NC, WI, TX, and FL---near-zero by construction of compactness-driven
redistricting that treats geographic space symmetrically. Enacted maps produce
substantially larger Partisan Bias ($-0.05$ to $-0.09$) reflecting deliberate
boundary choices that systematically favor Republican seat shares at the 50\%
vote-share pivot. The gap is robust to the choice of swing model: uniform swing,
heteroskedastic correction, and district-specific historical swing all reach the
same qualitative conclusion, with the heteroskedastic correction narrowing the
gap by at most 15\% for Texas.

The uniform swing limitation is real but manageable. Expert witnesses should
document it, apply the heteroskedastic correction as a robustness check, and
note that the correction narrows but does not reverse the qualitative conclusion.
For NC and WI, where safe seats are less prevalent, the correction is negligible
and the uniform swing result stands without qualification.

For state court expert witnesses, Partisan Bias is best presented alongside the
seats-votes curve (L.5), which provides the full picture of which the Bias is
a scalar summary. The combination---a graphical exhibit (the curve) and a
scalar number (the Bias)---provides courts with multiple ways to understand the
same underlying partisan asymmetry, reducing the risk that any single exhibit
is dismissed as insufficient.

\bigskip

\noindent\textit{The author thanks the redistricting research team for data
access and algorithmic development. All errors are the author's own.}
